\chapter{Bilan et ouverture : le schéma TFHE}

\section{Point historique et inventeurs}

Le schéma \textbf{TFHE} (\textit{Fast Fully Homomorphic Encryption over the Torus}) constitue l'une des avancées les plus significatives de la dernière décennie en cryptographie. Introduit en 2016, il est le fruit des travaux d'une équipe de chercheurs composée de \textbf{Ilaria Chillotti, Nicolas Gama, Mariya Georgieva et Malika Izabachène}. 

Ce schéma s'inscrit comme une amélioration radicale du schéma FHEW (2014). Son objectif principal était de lever le verrou technologique du \textit{bootstrapping}, en le rendant suffisamment rapide (de l'ordre de la milliseconde) pour permettre des calculs complexes et arbitrairement profonds sur des données chiffrées en temps quasi-réel.

\section{Comparaison entre les schémas DGHV et CKKS}

Le passage du DGHV au CKKS illustre l'évolution des priorités : de la preuve de concept arithmétique vers la performance applicative.

\begin{table}[h!]
\centering
\renewcommand{\arraystretch}{1.5}
\begin{tabular}{|l|p{5cm}|p{5cm}|}
\hline
\textbf{Caractéristique} & \textbf{DGHV} & \textbf{CKKS} \\ \hline
\textbf{Espace de calcul} & Entiers relatifs ($\mathbb{Z}$) & Nombres complexes / réels ($\mathbb{C}$ ou $\mathbb{R}$) \\ \hline
\textbf{Type de calcul} & Arithmétique exacte & Arithmétique approchée \\ \hline
\textbf{Avantage majeur} & Simplicité conceptuelle reposant sur les entiers. & Traitement de vecteurs (SIMD) pour le \textit{Machine Learning}. \\ \hline
\textbf{Limites} & Taille des chiffrés trop importante en pratique. & Inadapté aux calculs logiques exacts. \\ \hline
\end{tabular}
\caption{Comparaison synthétique des schémas DGHV et CKKS}
\end{table}

Alors que DGHV et CKKS se heurtent au coût du rafraîchissement du bruit, le TFHE propose une troisième voie : le calcul binaire exact et ultra-rapide sur une structure continue.

\section{Fondements mathématiques}

\subsection{Le tore}

\begin{definition}[Tore]
On appelle \textbf{tore} (en dimension 1) le quotient $\mathbb{T} = \mathbb{R} / \mathbb{Z}$.
\end{definition}

\begin{proposition}
Chaque élément de $\mathbb{T}$ admet un représentant unique dans l’intervalle $[0,1[$.
\end{proposition}

\subsection{Le problème TLWE (Torus Learning With Errors)}

La sécurité du TFHE repose sur le problème TLWE.

\begin{definition}[Échantillon TLWE]
Soit $k \ge 1$ et $\mathbf{s} \in \{0,1\}^k$ une clé secrète binaire. Un échantillon TLWE d'un message $\mu \in \mathbb{T}$ est un couple $(\mathbf{a}, b) \in \mathbb{T}^k \times \mathbb{T}$ tel que :
\[
b = \sum_{i=1}^{k} a_i \cdot s_i + e + \mu \pmod 1
\]
où $\mathbf{a}$ est un vecteur aléatoire uniforme sur $\mathbb{T}^k$ et $e$ est un bruit (c'est une gaussienne de faible amplitude).
\end{definition}

\subsection{Encodage des bits}

Pour maximiser la résistance au bruit, on encode généralement un bit $m \in \{0,1\}$ par $\mu = m \cdot \frac{1}{4} \in \mathbb{T}$. Ainsi, le bit $0$ correspond à la valeur $0$ et le bit $1$ à la valeur $0.25$.

\subsection{Fonctions booléennes}

TFHE permet d’évaluer des portes logiques (AND, OR, XOR) directement sur des échantillons TLWE. Comme toute fonction booléenne peut être décomposée en ces portes, TFHE est capable d'exécuter n'importe quel circuit logique de manière chiffrée. C'est à dire que n'importe quelle opération logique complexe peut être réduite à une suite de petites opérations élémentaires (en somme, multiplication...)

\section{Gestion du bruit et rafraîchissement}

\subsection{Le Bootstrapping Programmable}

Lors d’opérations successives, le bruit $e$ s'accumule. Si l'amplitude du bruit dépasse $1/8$, la distinction entre $0$ et $1/4$ devient impossible, rendant le déchiffrement erroné.

Le bootstrapping programmable est l'innovation majeure du schéma TFHE qui transforme une étape de maintenance technique en une étape de calcul utile. Voici un résumé de ses caractéristiques essentielles :

\begin{itemize}
    \item \textbf{Double fonction :} Contrairement au bootstrapping classique qui ne fait que \guillemotleft~nettoyer~\guillemotright\ le bruit d'un chiffré, le bootstrapping programmable réduit le bruit tout en évaluant simultanément une fonction logique ou mathématique.
    \item \textbf{Utilisation de LUT (Look-Up Table) :} Le processus s'appuie sur une table de correspondance. Pendant l'opération de rafraîchissement, le système applique la fonction encodée dans cette table directement à la donnée chiffrée.
    \item \textbf{Calcul non-linéaire :} Il permet de calculer efficacement des fonctions complexes ou non-linéaires (comme des comparaisons, des seuils ou des fonctions d'activation) qui sont d'ordinaire très coûteuses en chiffrement homomorphe.
    \item \textbf{Vitesse et continuité :} Grâce à l'arithmétique sur le tore, cette opération s'effectue en quelques millisecondes. Cela permet d'enchaîner un nombre virtuellement infini d'opérations puisque le niveau de bruit est réinitialisé après chaque calcul.
\end{itemize}

\section{Conclusion}

En conclusion, le schéma TFHE marque un tournant décisif dans l'évolution du chiffrement homomorphe. Si le schéma DGHV a posé les bases théoriques sur les entiers et CKKS a optimisé le calcul numérique approché, le TFHE apporte la brique manquante : l'efficacité du calcul logique exact.

Grâce à l'utilisation ingénieuse du tore réel et à l'introduction du bootstrapping programmable, ce schéma transforme une contrainte mathématique — le bruit — en un levier de calcul puissant. Cette capacité à traiter des fonctions booléennes en quelques millisecondes permet de s'affranchir des limites de profondeur des circuits, offrant ainsi une flexibilité d'utilisation jusqu'alors inédite.

Le TFHE ne se contente pas de sécuriser les données au repos ou en transit ; il instaure un \textbf{modèle de protection active}. En permettant l'exécution de programmes complets sur des données chiffrées sans compromis sur la précision, il s'impose comme une solution de référence pour les architectures de demain.